\documentclass{article} % For LaTeX2e
\usepackage{iclr2027_conference,times}

% Optional math commands from https://github.com/goodfeli/dlbook_notation.
\input{math_commands.tex}

\usepackage{hyperref}
\usepackage{url}

% --- added for this draft -------------------------------------------------
\usepackage{amsmath,amssymb,amsthm}
\usepackage{booktabs}
\usepackage{enumitem}

\theoremstyle{plain}
\newtheorem{theorem}{Theorem}
\newtheorem{proposition}[theorem]{Proposition}
\newtheorem{lemma}[theorem]{Lemma}
\theoremstyle{definition}
\newtheorem{definition}[theorem]{Definition}
\theoremstyle{remark}
\newtheorem{remark}[theorem]{Remark}
% -------------------------------------------------------------------------

\title{Diagnosing the implicit bias of gradient descent at the operator level}

\author{}

\newcommand{\fix}{\marginpar{FIX}}
\newcommand{\new}{\marginpar{NEW}}

%\iclrfinalcopy % Uncomment for camera-ready version, but NOT for submission.
\begin{document}


\maketitle

\begin{abstract}
\end{abstract}

% =========================================================================
\section{Introduction}
\label{sec:intro}

A deep network trained by gradient descent does not merely reach \emph{a} minimiser of its
training loss; on an over-parameterised model it reaches a particular one, and which one it
reaches is a property of the optimiser rather than of the objective. This selection effect is
usually called the \emph{implicit bias} of the optimiser, and for factored linear models it is by
now well characterised: depth drives the solution toward low rank, and the effect strengthens as
layers are added \citep{gunasekar2017implicit,arora2019implicit,li2021towards}.

Almost all of that characterisation is obtained by looking at \emph{endpoints}. One trains to
convergence, inspects the solution, and repeats the experiment with a different depth, width, or
initialisation; the optimiser's preference is then inferred from how the endpoint moves. This is
an indirect instrument, for a reason that is structural rather than practical: the comparison has
no reference point. To say that gradient descent's step is biased, one would have to say what an
unbiased step would have been, and in weight space there is no canonical answer --- the loss
supplies a gradient, and the gradient is precisely the object whose bias is in question.

This paper supplies the missing reference by changing the space in which the question is asked.
For a bias-free network built from linear maps and positively homogeneous gates, the function
computed at an input $x$ factors through a matrix,
\begin{equation}
  f(x) = P(x)\,x ,
  \label{eq:fx}
\end{equation}
where $P(x)$ is the input--output Jacobian at $x$ (defined precisely in
Section~\ref{sec:setup}). In this \emph{operator} view the loss has a direct opinion about what
should happen: it asks the operator to move along $-\partial \ell/\partial P(x)$. That request is
well defined, it does not mention the parameterisation, and every optimiser can be measured
against it.

What a weight update actually does to the operator is a different matter. The update $\Delta W$ is
applied to the factors, and the factorisation converts it into an operator change through a linear
map $M$ that depends on the current weights. Two questions then become separable, and both are
measurable:

\begin{enumerate}[leftmargin=*,itemsep=2pt,topsep=3pt]
  \item \textbf{How much of what the loss asked for is achievable at all?} The range of $M$ may
        not contain the requested operator change. This is a property of the architecture and the
        current weights, and \emph{no} optimiser can do better than it.
  \item \textbf{Of the achievable part, how much does a given optimiser take?} Gradient descent,
        Adam, or any other rule produces some $\Delta W$; the operator change it induces can be
        compared in direction and magnitude with the best achievable one. This is a property of
        the update rule.
\end{enumerate}

Separating these two is the contribution of this paper. It requires being able to compute, at any
point in training, the weight update that realises a prescribed operator change as exactly as the
architecture permits and with the smallest possible norm --- what we call the \emph{reference
step}. We give solvers for it in two regimes: exactly, in closed form, for deep linear chains
(Section~\ref{sec:linear}), and as a least-squares problem on the step map for gated networks
(Section~\ref{sec:nonlinear}).

\paragraph{Contributions.}
\begin{enumerate}[label=\textbf{C\arabic*.},leftmargin=*,itemsep=3pt,topsep=3pt]
  \item \textbf{A reference step targeting the operator.} Given a prescribed change of the
        operator, we compute the minimum-norm weight update realising it, up to the reachable set.
        The construction is a preconditioned gradient step (Remark~\ref{rem:precond}) and we do not
        claim the family as new: on a deep linear network with whitened inputs it reduces exactly to
        the natural gradient of \citet{bernacchia2018exact}. What is new is the target. Prescribing
        an operator change constrains the full Jacobian at every input, where prescribing an output
        change --- as in the closest prior formulation, \citet{evans2026theory} --- constrains only
        its action along that input; the two separate precisely when the network is gated or the
        inputs are not whitened (Sections~\ref{sec:setup}, \ref{sec:nonlinear}).
  \item \textbf{An exact solution of the prescribed-step problem at depth.} Realising a target
        operator exactly through $L$ factors, with minimal total norm, is a different problem from
        the linearised step above, and we solve it in closed form for deep linear chains
        (Section~\ref{sec:s1}). \citet{evans2026theory} solve the single-layer case exactly and
        bound the multi-layer case by Lipschitz arguments.
  \item \textbf{A decomposition of the bias into an architectural and an algorithmic part.}
        Reachability measures what the architecture forbids; deflection measures what the
        optimiser gives up within what is allowed. The first is optimiser-independent, the second
        is scale-invariant and therefore independent of the learning rate
        (Section~\ref{sec:diagnostics}).
  \item \textbf{A solver that is exact from degenerate points.} The prescribed-step problem admits
        an exact recursive solution even at parameter values where every derivative-based method
        is stationary, such as an all-zero chain or a deep chain whose product has underflowed
        (Section~\ref{sec:s1}).
  \item \textbf{Balancedness as a property of minimal realisations.} The layer-wise balance
        condition that underlies the classical deep linear analysis follows from minimality of the
        realisation alone, without invoking gradient flow (Section~\ref{sec:conserve}).
  \item \textbf{Measurements.} What the decomposition finds across model classes,
        initialisations, and optimisers (Sections~\ref{sec:linear-exp}, \ref{sec:nonlinear-exp}).
\end{enumerate}

% =========================================================================
\section{Setup: factored models, the operator, and the reference step}
\label{sec:setup}

This section defines every object used later, in order of dependence. Nothing here is specific to
a particular nonlinearity: the same definitions serve the deep linear networks of
Section~\ref{sec:linear} and the gated networks of Section~\ref{sec:nonlinear}, which is what
makes the diagnostics of Section~\ref{sec:diagnostics} comparable across model classes.

\subsection{Factored models}
\label{sec:models}

Throughout, a model has $L$ weight matrices $W_\ell \in \mathbb{R}^{n_\ell \times n_{\ell-1}}$ for
$\ell = 1,\dots,L$, and no bias terms. We write $W = (W_1,\dots,W_L)$ for the collection and
$\Delta W = (\Delta W_1,\dots,\Delta W_L)$ for an update. Between consecutive weight matrices sits
a \emph{gate} $D_\ell$, a matrix that may depend on the input. We consider four instances.

\begin{description}[leftmargin=*,itemsep=2pt,topsep=3pt]
  \item[Deep linear network (DLN).] $D_\ell = I$ for all $\ell$.
  \item[Fixed-gate linear network (FGLN).] $D_\ell$ is a fixed diagonal matrix, the same for every
        input. This interpolates between the linear and the gated case: the contexts below become
        input-dependent in form but not in value.
  \item[ReLU multilayer perceptron.] $D_\ell(x) = \mathrm{diag}\big(\mathbf{1}[z_\ell(x) > 0]\big)$,
        where $z_\ell(x)$ is the pre-activation at layer $\ell$.
  \item[CReLU multilayer perceptron.] The concatenated rectifier doubles the feature dimension,
        $c = [\,\mathrm{relu}(z)\,;\,\mathrm{relu}(-z)\,]$, which can be written $c = D(z)\,z$ with
        the rectangular gate
        \begin{equation}
          D(z) \;=\; \begin{bmatrix} \mathrm{diag}(\mathbf{1}[z>0]) \\[2pt]
                                     -\,\mathrm{diag}(\mathbf{1}[z<0]) \end{bmatrix}
          \in \mathbb{R}^{2d \times d},
          \qquad D(z)^\top D(z) = I .
          \label{eq:crelu}
        \end{equation}
        The isometry $D^\top D = I$ is what makes this architecture useful to us, for reasons
        given in Section~\ref{sec:lookslinear}.
\end{description}

All four gates are positively homogeneous of degree zero in the pre-activation, so the network is
positively homogeneous of degree one and \eqref{eq:fx} holds exactly rather than approximately.

\subsection{The operator and its contexts}
\label{sec:contexts}

\begin{definition}[Operator]
The \emph{operator} at input $x$ is
\begin{equation}
  P(x) \;=\; W_L\, D_{L-1}(x)\, W_{L-1} \cdots D_1(x)\, W_1
  \;\in\; \mathbb{R}^{n_L \times n_0},
  \label{eq:P}
\end{equation}
so that $f(x) = P(x)\,x$. Where the gates are input-independent (DLN, FGLN) we write $P$ without
the argument, and $P$ is then shared by all inputs.
\end{definition}

For a gated network $P(x)$ is the Jacobian $\partial f/\partial x$ almost everywhere, so
\eqref{eq:P} is not a new object but the familiar one, written in a form that exposes how each
layer enters.

\begin{definition}[Contexts]
The \emph{right contexts} and \emph{left contexts} are defined by the recursions
\begin{equation}
  B_1 = I, \qquad B_{\ell+1} = D_\ell W_\ell B_\ell,
  \qquad\qquad
  A_L = I, \qquad A_\ell = A_{\ell+1} W_{\ell+1} D_\ell .
  \label{eq:contexts}
\end{equation}
\end{definition}

They are exactly the parts of the chain to the right and to the left of a given layer, and by
construction
\begin{equation}
  P(x) \;=\; A_\ell(x)\, W_\ell\, B_\ell(x)
  \qquad \text{for every } \ell = 1,\dots,L .
  \label{eq:AWB}
\end{equation}
Equation~\eqref{eq:AWB} is the only structural fact we use, and it holds verbatim for all four
model classes of Section~\ref{sec:models}; only the value of $D_\ell$ changes. This is why a
single implementation and a single set of diagnostics cover the whole family.

\subsection{What the loss asks of the operator}
\label{sec:target}

Let $\ell(\cdot)$ be a per-example loss and $\mathcal{L}$ the average over a batch. Because
$f(x) = P(x)x$, the loss can be differentiated with respect to the operator directly.

\begin{definition}[Operator gradient]
$G(x) \;:=\; \partial \ell / \partial P(x) \in \mathbb{R}^{n_L \times n_0}$.
\end{definition}

For the two losses used in this paper $G$ is rank one:
\begin{equation}
  G(x) = \big(f(x) - y\big)\,x^\top \quad\text{(squared loss)},
  \qquad
  G(x) = \big(\mathrm{softmax}(f(x)) - e_y\big)\,x^\top \quad\text{(cross-entropy)} .
  \label{eq:G}
\end{equation}
Combining \eqref{eq:AWB} with the chain rule gives the familiar backpropagated gradient in the
notation above,
\begin{equation}
  \frac{\partial \mathcal{L}}{\partial W_\ell}
  \;=\; \sum_x A_\ell(x)^\top\, G(x)\, B_\ell(x)^\top .
  \label{eq:backprop}
\end{equation}

\begin{definition}[Target]
Given a step size $\eta > 0$, the \emph{target} is the collection of operator changes the loss
requests, one per input:
\begin{equation}
  \Delta_\star \;:=\; \big(-\eta\, G(x)\big)_x .
  \label{eq:target}
\end{equation}
\end{definition}

The target is defined without reference to the parameterisation. It is the statement ``move the
operator this way'', and it is the reference against which everything below is measured.

\subsection{The step map and the reference step}
\label{sec:stepmap}

An update $\Delta W$ changes the operator. To first order, substituting into \eqref{eq:P} and
collecting the terms linear in $\Delta W$,
\begin{equation}
  P\big(x;\, W + \Delta W\big) - P\big(x;\, W\big)
  \;=\; \sum_{\ell=1}^{L} A_\ell(x)\, \Delta W_\ell\, B_\ell(x) \;+\; O\!\left(\|\Delta W\|^2\right).
  \label{eq:firstorder}
\end{equation}

\begin{definition}[Step map]
The \emph{step map} is the linear operator sending a weight update to the collection of induced
operator changes,
\begin{equation}
  M(\Delta W) \;:=\; \Big( \textstyle\sum_{\ell} A_\ell(x)\, \Delta W_\ell\, B_\ell(x) \Big)_x ,
  \label{eq:M}
\end{equation}
with adjoint $\;(M^\top H)_\ell = \sum_x A_\ell(x)^\top H_x B_\ell(x)^\top$. We write $\Pi$ for the
orthogonal projector onto $\mathrm{range}(M)$.
\end{definition}

Comparing the adjoint with \eqref{eq:backprop} shows that gradient descent's weight step is
$\Delta W^{\mathrm{GD}} \propto M^\top \Delta_\star$. It is the adjoint applied to the target, not
the solution of any equation involving the target.

The gated architectures do not, in general, admit every requested operator change: $\Delta_\star$
need not lie in $\mathrm{range}(M)$. The best available request is therefore its projection, and
the reference step is the smallest update realising it.

\begin{definition}[Reference step]
\label{def:ref}
The \emph{reference step} is
\begin{equation}
  \Delta W^{\mathrm{ref}}
  \;=\; \arg\min \; \|\Delta W\|_F
  \quad \text{subject to} \quad
  M(\Delta W) = \Pi\,\Delta_\star ,
  \qquad\text{equivalently}\qquad
  \Delta W^{\mathrm{ref}} = M^{+} \Delta_\star ,
  \label{eq:ref}
\end{equation}
where $M^{+}$ is the Moore--Penrose pseudo-inverse and $\|\cdot\|_F$ is the Euclidean norm on the
concatenated update.
\end{definition}

Minimum norm is not an arbitrary tie-break. The constraint in \eqref{eq:ref} fixes only the
first-order effect on the operator; any element of $\ker M$ may be added without changing it, and
such an addition moves the weights while leaving the function unchanged to first order. Requiring
the smallest update is what makes the reference a statement about the operator alone. It also has
a consequence for the classical theory, developed in Section~\ref{sec:conserve}.

\begin{remark}[The reference step is a preconditioned gradient step]
\label{rem:precond}
Since $M^{+} = (M^\top M)^{+} M^\top$ and $M^\top \Delta_\star = -\eta\,\nabla_W \mathcal{L}$ by
\eqref{eq:backprop}, Definition~\ref{def:ref} is equivalent to
\begin{equation}
  \Delta W^{\mathrm{ref}} \;=\; -\eta\,\big(M^\top M\big)^{+}\,\nabla_W \mathcal{L} .
  \label{eq:precond}
\end{equation}
The reference step is therefore a preconditioned gradient step, with the preconditioner given by
the pseudo-inverse of the Gram of the operator Jacobian. It belongs to the Gauss--Newton and
natural-gradient family, and we state this plainly rather than presenting it as a separate
construction. What distinguishes it is \emph{which} metric is pulled back. Gauss--Newton uses the
Jacobian of the \emph{outputs}; \eqref{eq:precond} uses the Jacobian of the \emph{operator}. These
coincide exactly for a deep linear network on whitened inputs --- there
$\sum_i x_i x_i^\top \propto I$ makes the two Gram matrices proportional, and the reference step
reduces to the exact natural gradient of \citet{bernacchia2018exact} --- and separate otherwise.
Measured on a depth-3 network, the cosine between the two steps is $1.000$ for deep linear with
whitened inputs, $0.881$ for deep linear with raw inputs, and $0.690$ and $0.388$ for ReLU with and
without whitening. The separation has a clear cause: prescribing $\Delta P(x)$ constrains the whole
Jacobian at each input, whereas prescribing $\Delta f(x)$ constrains only its action along $x$.
\end{remark}

\subsection{What we measure}
\label{sec:diagnostics}

\begin{definition}[Reachability]
\begin{equation}
  \alpha \;:=\; \frac{\|\Pi \Delta_\star\|}{\|\Delta_\star\|} \;\in\; [0,1].
\end{equation}
\end{definition}
$\alpha$ is the fraction of the requested operator change that the architecture admits at the
current weights. It depends on the model and the point, \emph{not} on the optimiser: no update
rule can realise more than $\alpha$. For a deep linear network with a single shared target,
$\alpha = 1$ generically; for gated networks $\alpha < 1$, and Section~\ref{sec:nonlinear-exp}
measures it.

\begin{definition}[Deflection]
Let $\Delta W_m$ be the update some optimiser would take from the current weights. We compare it
with the reference in each of the two spaces:
\begin{equation}
  \underbrace{\cos_{\mathrm{op}} \;=\; \cos\!\big(\,\overbrace{M(\Delta W_m)}^{\text{operator change it causes}},\
      \overbrace{\Pi \Delta_\star}^{\text{best available}}\,\big)}_{\text{operator space}},
  \qquad
  \underbrace{\cos_{\mathrm{w}} \;=\; \cos\!\big(\,\overbrace{\Delta W_m}^{\text{step it takes}},\
      \overbrace{\Delta W^{\mathrm{ref}}}^{\text{step it should have taken}}\,\big)}_{\text{weight space}},
\end{equation}
together with the corresponding norm ratios $\|M(\Delta W_m)\|/\|\Pi\Delta_\star\|$ and
$\|\Delta W_m\|/\|\Delta W^{\mathrm{ref}}\|$.
\end{definition}

In words, there are only three objects and each is the answer to a plain question.
$\Delta_\star$ is \emph{what the loss asked the operator to do}. $\Pi\Delta_\star$ is \emph{how much
of that the architecture can actually deliver} --- and when the model is over-parameterised
relative to the batch, $M$ is surjective, $\Pi\Delta_\star=\Delta_\star$, $\alpha=1$, and this
distinction disappears entirely. $\Delta W^{\mathrm{ref}}$ is \emph{the weight update that delivers
it}, at least norm.

\textbf{Why both spaces.} The two cosines can disagree, and the disagreement is exactly what the
reference makes visible. Any element of $\ker M$ may be added to a weight update without changing
the operator at all, so a high $\cos_{\mathrm{op}}$ beside a low $\cos_{\mathrm{w}}$ says an
optimiser reached roughly the right operator motion \emph{by a different route through weight
space}. Asking that question requires a reference that is itself a weight update, which is what
distinguishes this construction from a purely functional comparison.

Both are cosines, hence invariant to the scale of $\Delta W_m$ and therefore \textbf{independent of
the learning rate}; only the norm ratios carry one, which is why they are reported separately. This
matters for the evidential status of the comparison: a claim expressed in these quantities cannot
be answered by asserting that a baseline was mis-tuned.

\begin{remark}[Fidelity is not quality]
\label{rem:fidelity}
$\Delta_\star = -\eta\,\partial\mathcal{L}/\partial P$ is steepest descent \emph{in operator
space}. On an anisotropic loss, steepest descent is itself an indirect route to the minimiser: it
removes the steep directions first and then advances slowly along the flat ones. Consequently an
optimiser may score \emph{lower} on $\cos_{\mathrm{op}}$ and still reach the optimum \emph{sooner}
--- a method whose step approximates a whitened or second-order direction will do exactly that.
$\cos_{\mathrm{op}}$ quantifies how far the factorisation deflects the step the loss asked for. It
is not a ranking of optimisers and must not be read as one.
\end{remark}

\begin{remark}[$\eta$ is an operator step, not a learning rate]
The target \eqref{eq:target} asks the operator to move by $\eta\|G\|$, so $\eta$ reads directly as
a fraction of $\|P(x)\|$ and is comparable across depths and architectures. A learning rate is not:
the same numerical value means different things at different depths because the map from weights
to operator changes with depth.
\end{remark}

\begin{remark}[Gradient descent as a truncation]
\label{rem:krylov}
Solving \eqref{eq:ref} with a Krylov least-squares method started from zero, the first iterate is
proportional to $M^\top \Delta_\star$, which by \eqref{eq:backprop} is exactly the gradient descent
direction. The reference step is what the same solve returns on convergence. This is a convenient
implementation fact --- it means one routine produces both arms, and a truncation parameter
interpolates between them --- and we use it in Section~\ref{sec:s2}. We draw no further conclusion
from it.
\end{remark}

% =========================================================================
\section{What is already known}
\label{sec:sota}

\subsection{Dynamics and implicit bias in deep linear models}

The deep linear network is the standard setting in which implicit bias can be analysed rather than
only observed, because the function computed is a single matrix while the optimisation remains
non-convex. \citet{saxe2014exact} obtained exact solutions of the learning dynamics under gradient
flow from balanced and aligned initialisation with whitened inputs, exhibiting the sequential,
mode-by-mode learning that has since organised much of the field.
\citet{arora2018optimization} showed that over-parameterisation by depth acts as a preconditioner:
under gradient flow with a balanced initialisation, $W_{j+1}^\top W_{j+1} = W_j W_j^\top$, the
induced motion of the end-to-end matrix is
\begin{equation}
  \dot P \;=\; -\sum_{j=1}^{L}\big(PP^\top\big)^{\frac{L-j}{L}}\,G\,\big(P^\top P\big)^{\frac{j-1}{L}} ,
  \label{eq:arora}
\end{equation}
which on the singular values of $P$ reduces to a per-mode gain $L\,s^{2-2/L}$. Both hypotheses are
essential and neither is generic: \eqref{eq:arora} is a statement about gradient \emph{flow} from a
\emph{balanced} initialisation.

\citet{du2018algorithmic} showed that the layer-wise imbalance
$W_{\ell+1}^\top W_{\ell+1} - W_\ell W_\ell^\top$ is conserved under gradient flow for homogeneous
models, which is what makes balancedness a persistent property rather than an initial one, and
extends beyond the linear case. \citet{bah2022learning} identified the factor flow with a
Riemannian gradient flow on the manifold of fixed-rank matrices, and \citet{ablin2020deep} showed
that a deep \emph{orthogonally constrained} linear network is equivalent to a single layer, which
delimits how much of the depth effect survives when the factors are restricted.
\citet{wendin2026gradient} survey the gradient-flow equations for these models, and
\citet{cohen2024lectures} give a consolidated account of the area together with a list of
questions that remain open --- among them the behaviour of unbalanced initialisation beyond depth
two, which is the regime the reference step of Section~\ref{sec:stepmap} is designed to reach.

On the bias itself, \citet{gunasekar2017implicit} conjectured that gradient descent on a two-layer
factorisation with small initialisation and small steps converges to the minimum nuclear norm
solution, connecting the dynamics to the classical convex surrogate for rank
\citep{srebro2005maximum,recht2010guaranteed}. \citet{arora2019implicit} showed that depth
strengthens the tendency toward low rank in deep matrix factorisation, and
\citet{razin2020implicit} then showed that no norm --- nuclear or otherwise --- can explain the
bias in general, which redirected the field from characterising the endpoint by a penalty toward
characterising the trajectory. \citet{li2021towards} showed that the flow from infinitesimal
initialisation performs a greedy low-rank search, and \citet{jacot2021saddle} described the
saddle-to-saddle dynamics in which the effective rank increases by one at each saddle escape.
\citet{tarmoun2021understanding}, \citet{min2021explicit} and \citet{kunin2024get} analysed how the
initialisation --- in particular its imbalance --- controls both the convergence rate and the
resulting bias, showing that unbalanced initialisation is not a perturbation of the balanced theory
but a qualitatively different regime.

\subsection{Beyond the linear case}

For nonlinear networks the picture is less complete. \citet{timor2023implicit} established
regularisation toward rank minimisation in ReLU networks under stated conditions, and
\citet{jacot2023implicit} introduced a notion of rank for nonlinear functions under which large
depth induces a bias toward low rank. \citet{bantzis2025saddle} extended the saddle-escape analysis
to deep ReLU networks, finding a low-rank bias at the first escape, and \citet{rawal2026saddle} give
a theory of saddle escape in deep nonlinear networks. \citet{gatedlinear2022implicit} analysed
generalized gated linear networks, which is the closest published setting to the FGLN model of
Section~\ref{sec:models}. \citet{huh2021lowrank} documented the low-rank simplicity bias
empirically across architectures. A separate line studies the average gradient outer product as the
mechanism of feature learning \citep{radhakrishnan2024mechanism,beaglehole2024agop}.

Two further results bear directly on the objects we measure. \citet{pennington2017resurrecting}
showed that ReLU networks cannot achieve dynamical isometry --- the input--output Jacobian cannot
have a well-conditioned spectrum at depth --- which is why the reachability $\alpha$ of
Section~\ref{sec:diagnostics} cannot be expected to equal one for ReLU, and why a nonlinear
architecture that \emph{is} conditioned by construction (Section~\ref{sec:lookslinear}) is needed as
a control. \citet{wang2016extended} proposed the extended data Jacobian matrix as a diagnostic,
which is an operator-level measurement in our sense, though used for model comparison rather than
for scoring update rules. \citet{haas2026why} analyse how the Jacobian spectra of deep networks
separate with depth, and establish the conservation of the gate-absorbed imbalance that
Section~\ref{sec:conserve} recovers from minimality.

\subsection{Geometry-aware and second-order optimisation}

The reference step of Definition~\ref{def:ref} is a preconditioned gradient step
(Remark~\ref{rem:precond}), so it belongs squarely to this literature and we position it there
rather than apart from it. \citet{bernacchia2018exact} computed the exact natural gradient for deep
linear networks and observed that it annihilates precisely the depth-induced preconditioner in
\eqref{eq:arora}. That is the closest relation, and it is an identity rather than an analogy: on a
deep linear network with whitened inputs the reference step \emph{is} their natural gradient, since
the operator and output Gram matrices are then proportional. It separates from Gauss--Newton and
from natural gradient once the inputs are not whitened or the gates are input-dependent, because
prescribing $\Delta P(x)$ constrains the entire Jacobian at $x$ while prescribing $\Delta f(x)$
constrains only its action along $x$.

\citet{evans2026theory} give the closest published formulation of the underlying problem: they
derive minimal-norm weight perturbations achieving a specified change in \emph{output}, exactly for
a single layer, and compare against Lipschitz-based bounds for the multi-layer case; their purpose
is robustness certification rather than optimiser analysis. The two differences from our setting are
the target --- output versus operator --- and the treatment of depth, where we solve the multi-layer
prescribed-step problem exactly (Section~\ref{sec:s1}). \citet{frerix2018proximal} recast
backpropagation as a sequence of proximal steps, and \citet{dufortlabbe2026layerwise} develop a
layerwise LQR view of geometry-aware optimisation. The distinction of purpose is worth stating
once: all of these are optimisation methods proposed to train faster, whereas we use the algebra as
a \emph{measuring instrument}, and Remark~\ref{rem:fidelity} explains why the two goals must not be
conflated.

\subsection{The gap}

Across this literature the bias is inferred from endpoints, and its architectural and algorithmic
components are not separated. Statements of the form ``depth strengthens the low-rank bias'' are
established by varying depth and observing the endpoint; they do not distinguish a change in what
the architecture \emph{permits} from a change in what the optimiser \emph{does} within what is
permitted. The reference step makes that separation computable, and the two diagnostics of
Section~\ref{sec:diagnostics} report the two halves.

\begin{table}[t]
\caption{Results we rely on or measure against, with the hypotheses each actually requires.}
\label{tab:sota}
\begin{center}
\small
\begin{tabular}{p{4.0cm}p{1.6cm}p{5.0cm}p{2.2cm}}
\toprule
\textbf{Result} & \textbf{Model} & \textbf{Hypotheses required} & \textbf{Source} \\
\midrule
Exact dynamics; modes learned largest-first & DLN &
gradient flow; balanced \emph{and} aligned init; whitened inputs & \citet{saxe2014exact} \\
Induced step \eqref{eq:arora}; per-mode gain $L s^{2-2/L}$ & DLN &
gradient \emph{flow}; balanced init & \citet{arora2018optimization} \\
Layer imbalance conserved & homogeneous &
gradient flow; homogeneity & \citet{du2018algorithmic} \\
Factor flow $=$ Riemannian flow on fixed-rank manifold & DLN & --- & \citet{bah2022learning} \\
Deep orthogonal linear $\equiv$ one layer & DLN & orthogonality constraint & \citet{ablin2020deep} \\
Bias toward minimum nuclear norm & 2-layer & small init and step; \emph{conjecture} & \citet{gunasekar2017implicit} \\
Depth strengthens the low-rank bias & deep MF & matrix sensing / completion & \citet{arora2019implicit} \\
No norm explains the bias & deep MF & specific instances & \citet{razin2020implicit} \\
Flow $\equiv$ greedy low-rank learning & 2-layer & infinitesimal init & \citet{li2021towards} \\
Saddle-to-saddle; rank $+1$ per escape & DLN & small init & \citet{jacot2021saddle} \\
Rank minimisation in ReLU networks & ReLU & stated conditions & \citet{timor2023implicit} \\
Large depth induces low nonlinear rank & nonlinear & notion of rank defined therein & \citet{jacot2023implicit} \\
ReLU cannot achieve dynamical isometry & ReLU & random init, large width & \citet{pennington2017resurrecting} \\
Exact natural gradient; cancels the depth preconditioner & DLN & invertible Fisher & \citet{bernacchia2018exact} \\
\bottomrule
\end{tabular}
\end{center}
\end{table}

% =========================================================================
\section{Part I --- Deep linear networks}
\label{sec:linear}

In a deep linear network the gates are the identity, so $P = W_L \cdots W_1$ is a single matrix
shared by every input. The target reduces to one matrix, $\Delta_\star = -\eta\,
\partial\mathcal{L}/\partial P$, and the step map \eqref{eq:M} maps weight updates into the space
of $n_L \times n_0$ matrices. This is the regime in which the reference step can be computed
exactly and in closed form, and in which the classical theory of Section~\ref{sec:sota} gives us
something to validate against.

\subsection{The prescribed-step problem}
\label{sec:prescribed}

Rather than solving \eqref{eq:ref} in its linearised form, in the linear case we can pose and
solve the exact, non-linearised problem.

\begin{definition}[Prescribed-step problem]
\label{def:prescribed}
Given a target matrix $P^\star$ and a depth $L$, find factors minimising the total squared norm
subject to realising the target exactly:
\begin{equation}
  \min_{W_1,\dots,W_L} \; \sum_{\ell=1}^{L} \|W_\ell\|_F^2
  \qquad \text{subject to} \qquad
  W_L W_{L-1} \cdots W_1 \;=\; P^\star .
  \label{eq:prescribed}
\end{equation}
\end{definition}

The solution is classical in form and is governed by a Schatten quasi-norm. Writing
$\|\cdot\|_{S_p}$ for the Schatten-$p$ quasi-norm, $\|A\|_{S_p} = (\sum_k \sigma_k(A)^p)^{1/p}$:

\begin{proposition}[Minimal realisation]
\label{prop:minimal}
The optimum of \eqref{eq:prescribed} equals $L \, \|P^\star\|_{S_{2/L}}^{2/L}$, and is attained by
any factorisation whose factors share the singular vectors of $P^\star$ in the chain and distribute
its singular values as $\sigma_k^{1/L}$ in every factor. In particular the optimal factorisation is
\emph{balanced}: $W_{\ell+1}^\top W_{\ell+1} = W_\ell W_\ell^\top$ for all $\ell$.
\end{proposition}

The proof is by a Bellman recursion over the depth with value function
$V_a(Y) = a\|Y\|_{S_{2/a}}^{2/a}$ for a sub-chain of depth $a$, together with the arithmetic
mean--geometric mean inequality to establish that the per-mode split is the equalising one; it is
given in Appendix~\ref{app:proofs}. The quasi-norm appearing here is the same object that arises
as a non-convex surrogate for rank, and whose direct minimisation is studied algorithmically by
\citet{shang2016scalable}; the difference is that it emerges for us as the value of
\eqref{eq:prescribed} rather than being imposed as a penalty.

Proposition~\ref{prop:minimal} already says something about the literature. Balancedness enters the
classical analysis as a hypothesis on the \emph{initialisation}, preserved thereafter by the
conservation law of \citet{du2018algorithmic}. Here it appears instead as a consequence of
minimality, with no dynamics involved. We return to this in Section~\ref{sec:conserve}.

\subsection{Degeneracy: where derivative-based methods stop and the prescribed step does not}
\label{sec:degeneracy}

The practically important feature of \eqref{eq:prescribed} is that it remains solvable at parameter
values where gradient information vanishes entirely.

\begin{definition}[Degenerate point]
A parameter value $W$ is \emph{degenerate} if $P = 0$ and all contexts vanish, i.e.\ $A_\ell = 0$
or $B_\ell = 0$ for every $\ell$.
\end{definition}

\begin{proposition}[Exactness from degeneracy]
\label{prop:degen}
At a degenerate point, $\partial \mathcal{L}/\partial W_\ell = 0$ for every $\ell$, so every
method whose update is a function of the gradient --- gradient descent, momentum, Adam, and any
preconditioned variant --- is stationary. The prescribed-step problem \eqref{eq:prescribed}
nevertheless has a solution for every $P^\star$ and every $L$, and the construction of
Section~\ref{sec:s1} attains it.
\end{proposition}

\begin{remark}[Why this is not a statement about the reference step]
\label{rem:twoproblems}
The distinction matters and is easy to lose. At a degenerate point $M = 0$, so the linearised
reference step \eqref{eq:precond} vanishes along with the gradient: as a preconditioned gradient
method it is stationary there too. Proposition~\ref{prop:degen} is a statement about the
\emph{exact} problem \eqref{eq:prescribed}, which asks for factors whose product equals a target
rather than for a small update whose first-order effect matches a requested change. That problem
does not reference the gradient at all, which is why it remains solvable where every gradient-based
method --- ours included --- stops.
\end{remark}

The all-zero chain is the extreme case, but the situation it models is not exotic. In a deep chain
under a standard random initialisation the product $W_L\cdots W_1$ has a spectrum that contracts
geometrically with depth; past a moderate $L$ the contexts are numerically indistinguishable from
zero in the directions that matter, and the gradient carries no usable signal even though it is not
formally zero. The degenerate point is the idealisation of that regime, and a construction that is
exact there is well defined throughout it.

\subsection{Minimality implies balancedness}
\label{sec:conserve}

The minimum-norm condition in Definition~\ref{def:ref} has a consequence that recovers, from a
different premise, the structure the classical analysis assumes.

\begin{proposition}[Conservation from minimality]
\label{prop:conserve}
Let $\Delta W$ solve \eqref{eq:ref}. Then there exists a single matrix $\Lambda$, the Lagrange
multiplier of the constraint, such that
\begin{equation}
  \Delta W_\ell \;=\; A_\ell^\top\, \Lambda\, B_\ell^\top
  \qquad \text{for every } \ell .
  \label{eq:kkt}
\end{equation}
Consequently $\Delta W_{\ell+1}^\top W_{\ell+1}$ and $\Delta W_\ell W_\ell^\top$ are transposes of
one another, and the imbalance $W_{\ell+1}^\top W_{\ell+1} - W_\ell W_\ell^\top$ is unchanged to
first order in the step.
\end{proposition}

The content of Proposition~\ref{prop:conserve} is that the shared multiplier in \eqref{eq:kkt} ---
one $\Lambda$ for all layers --- is forced by minimality, and the balance property follows from
that form alone. It requires neither gradient flow nor a balanced initialisation. Since gradient
descent's step $M^\top \Delta_\star$ is of the same form with $\Lambda = \Delta_\star$, the
classical conservation law is recovered as the special case in which the multiplier is the target
itself.

\subsection{Solver S1: exact minimal realisation by recursive splitting}
\label{sec:s1}

We solve \eqref{eq:prescribed} by recursion on the depth. Suppose the chain of depth $L$ is split
into a lower part of depth $a$ and an upper part of depth $b$, with $a + b = L$, so that
$P^\star = Y X$ where $X$ is realised by the lower $a$ factors and $Y$ by the upper $b$. Taking the
singular value decomposition $P^\star = U \Sigma V^\top$, we set
\begin{equation}
  Y_0 \;=\; U\, \Sigma^{\,t} , \qquad
  X_0 \;=\; \Sigma^{\,1-t}\, V^\top , \qquad
  t \;=\; \frac{b}{a+b} ,
  \label{eq:split}
\end{equation}
and recurse on $X_0$ with depth $a$ and on $Y_0$ with depth $b$. The exponent in \eqref{eq:split}
is \emph{depth-proportional}: each half receives the share of the singular values corresponding to
the number of factors it must still produce. An even split $t = 1/2$ is correct only when
$a = b$, and using it in general is what breaks exactness at odd depths.

The recursion bottoms out at $L = 1$, where the factor is the target itself. It never queries a
gradient, so it applies unchanged at a degenerate point, which establishes
Proposition~\ref{prop:degen}. The cost is that of one singular value decomposition per split,
$O(L)$ of them in total.

\begin{remark}[Relation to alternating least squares]
An obvious alternative is to solve for one factor at a time holding the others fixed, which is
block coordinate descent, or alternating least squares. That procedure converges to a realisation
of the target but not to the \emph{minimal} one: each sweep solves an exactly determined
sub-problem and has no reason to respect the global norm budget of
Proposition~\ref{prop:minimal}. We report the resulting gap in Appendix~\ref{app:als}.
\end{remark}

\subsection{Experiments}
\label{sec:linear-exp}

Each experiment below states its motivation, its exact configuration, the quantity measured with
its formula, its current status, the numbers obtained, and what may and may not be concluded from
them. Status is one of \textsc{done}, \textsc{partial} or \textsc{todo}.

\subsubsection{E1 --- Validating the machinery against a known closed form}

\paragraph{Motivation.}
Every later measurement is computed from the contexts \eqref{eq:contexts} and the step map
\eqref{eq:M}. Before any of them can be believed, the implementation must reproduce a result that
is known analytically. The per-mode gain of \citet{arora2018optimization} provides one, with the
advantage that it is exact rather than asymptotic, so agreement can be demanded to machine
precision rather than to within noise.

\paragraph{Setup.}
A balanced chain carrying a prescribed non-flat spectrum, $L \in \{2,4,8,16,32\}$, in both the deep
linear and the CReLU implementation (looks-linear blocks built from the same spectrum, so the gate
absorbs and $P$ remains non-flat). All diagnostics in float64 on CPU.

\paragraph{Measurement.}
With the transfer operator $T(H) := \sum_\ell A_\ell A_\ell^\top H\, B_\ell^\top B_\ell$, so that
the induced operator change is $-\eta\,T(G) + O(\eta^2)$, and the singular value decomposition
$P = \sum_k s_k u_k v_k^\top$, the \emph{mode gain} is
\begin{equation}
  c_k \;:=\; \big\langle u_k,\ T(u_k v_k^\top)\, v_k \big\rangle
      \;=\; \sum_{\ell} \big\| A_\ell^\top u_k \big\|^2 \, \big\| B_\ell v_k \big\|^2 .
  \label{eq:gain}
\end{equation}
We fit $\log c_k$ against $\log s_k$ and compare the slope with $2 - 2/L$ and the intercept
with $L$, which is what \eqref{eq:arora} predicts under balancedness.

\paragraph{Status.} \textsc{done}.

\paragraph{Results.}
The measured slope equals $2 - 2/L$ and the intercept equals $L$ to $10^{-14}$ for every depth
tested, in both implementations.

\paragraph{Conclusions.}
The contexts, the transfer operator and the gain measurement are correct as implemented, in the
gated as well as the ungated case. Two implementation traps are worth recording because each
produces a qualitatively wrong exponent rather than a noisy one. First, $c_k$ must be probed with
the rank-one direction $u_k v_k^\top$; probing with a random $H$ and forming
$\langle u_k, T(H) v_k\rangle / \langle u_k, H v_k\rangle$ equals $c_k$ only when the context Gram
matrices are diagonalised by the singular vectors of $P$, which holds under balancedness and fails
otherwise. Second, a balanced chain is not an orthogonal one: an orthogonal chain is balanced but
has a flat spectrum, from which no exponent is identifiable at all.

This experiment validates the instrument. It is not evidence about optimisers, and the gain
exponent is not used as a comparison metric anywhere in this paper: being a function of the current
weights alone, it takes the same value for two optimisers evaluated at the same point.

\subsubsection{E2 --- Operator-space trajectories}

\paragraph{Motivation.}
The diagnostics of Section~\ref{sec:diagnostics} are per-step scalars. This experiment shows the
same phenomenon as a picture, in a setting where the operator is two-dimensional so the trajectory
can be drawn exactly rather than projected, and where the loss is exactly quadratic so the geometry
is not in question.

\paragraph{Setup.}
A deep linear chain $W_L \cdots W_1 = P \in \mathbb{R}^{1 \times 2}$ on a two-feature regression:
$X \in \mathbb{R}^{256 \times 2}$ with feature scales $(1.0, 0.35)$ so the loss is anisotropic,
$y = Xw$ with $w$ drawn once per seed, hidden width $8$, $L \in \{2, 8, 32\}$, $5$ seeds,
$400$ steps, float64 on CPU. Arms: gradient descent and damped alternating least squares at
learning rate $0.06$, Adam at $0.015$ (Adam requires its own step size; at the gradient descent
value it oscillates, which would make the figure about mis-tuning rather than about geometry), and
the reference step at $\eta = 0.06$. The weight space has $1944$ parameters at $L = 32$; only the
operator is two-dimensional.

\paragraph{Measurement.}
At each step we record the realised operator change $\Delta P$ and its \emph{alignment} with what
the loss requested,
\begin{equation}
  \mathrm{align} \;:=\; \cos\!\Big(\Delta P,\ -\tfrac{\partial \mathcal{L}}{\partial P}\Big),
  \label{eq:align}
\end{equation}
together with the final distance $\|P - P^\star\|$ to the ordinary least squares solution. We
report the mean alignment, its minimum, and the fraction of steps on which it is negative --- that
is, on which the step moved the operator \emph{against} the direction the loss asked for.

\paragraph{Status.} \textsc{done}.

\paragraph{Results.}
Percentage of steps with negative alignment, mean over $5$ seeds with the per-seed range in
brackets. Runs that diverged are excluded and counted separately; a diverging run contributes an
alignment statistic computed from a handful of blown-up steps, which is not a property of the
method.

\begin{center}\small
\begin{tabular}{lccc}
\toprule
depth & gradient descent & Adam & reference \\
\midrule
$L=2$  & $0.3\%$ $[0,2]$ & $9.3\%$ $[0,22]$  & $0.0\%$ \\
$L=8$  & $1.0\%$ $[0,3]$ & $33.0\%$ $[28,37]$ & $0.0\%$ \\
$L=32$ & $0.0\%$ $[0,0]$ & $26.9\%$ $[15,37]$ & $0.1\%$ \\
\bottomrule
\end{tabular}
\end{center}

\paragraph{Conclusions.}
Adam moves the operator against the requested direction on a substantial fraction of steps at
depth, in every seed --- the per-seed range never approaches zero --- while gradient descent and the
reference essentially never do. The effect rises sharply from $L=2$ and then plateaus; with five
seeds the difference between $L=8$ and $L=32$ lies within the spread, so we claim a rise and a
plateau rather than a monotone law in depth.

Two cautions attach to this figure. It must be drawn in operator space: a model whose \emph{weight}
space is two-dimensional has a rank-one step map $M = [\,w_2\ \ w_1\,]$, whence $M^{+} \propto
M^\top$ and gradient descent coincides with the reference exactly, measured at
$\cos = 1.0000000000$. And alternating least squares requires ridge damping here; undamped, it is
anti-aligned on $100\%$ of steps at $L=2$ and diverges to $\|P - P^\star\| = 154$.

Finally, Remark~\ref{rem:fidelity} is visible in the numbers: at $L = 32$ Adam reaches
$\|P - P^\star\| = 0.0000$ while the reference reaches $0.1437$. Adam is the worse-aligned method
and the faster one. Alignment is fidelity to the operator gradient, not optimisation quality.

% =========================================================================
\section{Part II --- Gated networks}
\label{sec:nonlinear}

\subsection{What changes}
\label{sec:whatchanges}

Three things change at once when the gates become input-dependent, and it is worth separating
them.

\begin{enumerate}[leftmargin=*,itemsep=2pt,topsep=3pt]
  \item \textbf{The operator is no longer shared.} Each input has its own $P(x)$, and the target
        \eqref{eq:target} asks for a different operator change at each one. A single weight update
        must serve all of these requests simultaneously.
  \item \textbf{The requests are not jointly realisable.} Since one $\Delta W$ produces the whole
        collection $M(\Delta W)$, and the collection has far more degrees of freedom than
        $\Delta W$, in general $\Delta_\star \notin \mathrm{range}(M)$ and $\alpha < 1$. The
        reference step realises the projection $\Pi \Delta_\star$, which is the most any optimiser
        can do.
  \item \textbf{Each sample constrains only a rank-one slice of its own operator.} Fitting the
        data requires $P(x_i) x_i = y_i$, which fixes the action of $P(x_i)$ along $x_i$ and leaves
        the remaining directions free. This is why gated networks can exhibit genuine differences
        between optimisers at the same training loss, while a deep linear network with enough
        samples cannot: there the single shared $P$ is pinned by the data, and all methods that fit
        must agree.
\end{enumerate}

Point 3 is the reason the endpoint comparisons of Section~\ref{sec:nonlinear-exp} are posed for
gated networks and not for deep linear ones.

\subsection{CReLU and looks-linear initialisation: a conditioned nonlinear control}
\label{sec:lookslinear}

The ReLU network confounds two effects that we wish to separate: it is nonlinear, and at depth it
is ill-conditioned, necessarily so \citep{pennington2017resurrecting}. To vary gating and
conditioning independently we need a network that is genuinely nonlinear yet perfectly conditioned
by construction. The concatenated rectifier of \eqref{eq:crelu} provides one.

Write the layer weight in block form $W = [\,O \mid -O\,]$ with $O$ orthogonal --- the
\emph{looks-linear} initialisation. Then for every $z$,
\begin{equation}
  W D(z) \;=\; O\,\mathrm{diag}(\mathbf{1}[z>0]) \;+\; O\,\mathrm{diag}(\mathbf{1}[z<0]) \;=\; O ,
\end{equation}
because the two indicator masks are complementary. Hence $P(x) = O_L \cdots O_1$ is orthogonal, at
every input and at \emph{any} depth: all singular values equal one, and the operator is identical
across inputs at initialisation. The network is nonetheless a genuine CReLU network --- the gates
are real, they differ between inputs as soon as training moves the weights away from the
looks-linear point, and no reparameterisation makes it linear.

This gives the control the rest of Part II requires: CReLU with looks-linear initialisation is
nonlinear and perfectly conditioned; ReLU with a standard random initialisation is nonlinear and
ill-conditioned at depth; the deep linear network of Part I is linear. Differences attributed to
gating can then be distinguished from differences attributable to conditioning.

\subsection{Solver S2: least squares on the step map}
\label{sec:s2}

For gated networks we solve \eqref{eq:ref} directly as a least-squares problem,
\begin{equation}
  \Delta W^{\mathrm{ref}} \;=\; \arg\min_{\Delta W} \; \big\|M(\Delta W) - \Delta_\star\big\|^2 ,
  \qquad \text{minimum-norm solution},
\end{equation}
using LSQR \citep{paige1982lsqr}, which requires only products with $M$ and $M^\top$ and never
forms the matrix. By \eqref{eq:M} and its adjoint, both products are batched matrix
multiplications over layers and inputs, so the solve costs $O(k)$ applications of the forward and
backward context products for $k$ iterations. By Remark~\ref{rem:krylov}, $k = 1$ returns the
gradient descent direction and large $k$ returns the reference step, so a single routine produces
both arms.

\paragraph{A numerical caveat that must not be skipped.}
On these step maps, LSQR \emph{diverges after it has converged}: the Lanczos basis loses
orthogonality, and the iterate grows without bound while the least-squares residual remains at its
converged value. Running to a fixed large iteration count therefore returns a meaningless result
with no error raised. Two natural stopping tests do not detect this. The quantity $\bar\phi$
maintained internally by the LSQR recursion is monotonically non-increasing by construction and so
never signals anything. The residual $\|M(\Delta W) - \Delta_\star\|$ is blind to the component of
the failure that lies in $\ker M$, which is precisely where the iterate grows. The normal-equation
residual $\|M^\top(M(\Delta W) - \Delta_\star)\|$ detects both, and is the stopping rule we use; it
is not monotone, so the best iterate is retained rather than the last, and its attained value is
recorded alongside every measurement. We report quantitative diagnostics in
Appendix~\ref{app:numerics}.

\subsection{Experiments}
\label{sec:nonlinear-exp}

\subsubsection{E3 --- How much of the request is reachable at all}

\paragraph{Motivation.}
Reachability is the ceiling: it is a property of the architecture and the current weights, and no
optimiser can exceed it. It must therefore be measured before any optimiser is scored, since a low
score against a low ceiling means something different from a low score against a high one.

\paragraph{Setup.}
Teacher--student regression, input dimension $32$, $2000$ training and $1000$ test points, batch
$100$; student of depth $L = 8$ and width $32$; architectures $\{$deep linear, FGLN, ReLU,
CReLU$\}$ $\times$ initialisations $\{$Xavier, orthogonal$\}$, with looks-linear substituted for
orthogonal in the CReLU case. Probes every $100$ steps along a $1000$-step trajectory, $11$ probes
per cell, $8$ cells. Diagnostics in float64.

\paragraph{Measurement.}
$\alpha := \|\Pi \Delta_\star\| / \|\Delta_\star\| \in [0,1]$, with $\Pi \Delta_\star$ obtained
from the damped solve of Section~\ref{sec:s2}. We also record the solve quality
\begin{equation}
  \mathrm{ne} \;:=\; \frac{\big\|M^\top\big(M(\Delta W^{\mathrm{ref}}) - \Delta_\star\big)\big\|}
                          {\big\|M^\top \Delta_\star\big\|},
  \label{eq:ne}
\end{equation}
the relative normal-equation residual, at every probe, and discard any probe with
$\mathrm{ne} > 10^{-2}$. Since $\alpha$ is a projection ratio, $\alpha > 1$ is impossible and
serves as an independent check on the solve.

\paragraph{Status.} \textsc{done}. $44/44$ probes valid after the fix described below.

\paragraph{Results.}
$\alpha$ ranges from $0.089$ to $0.438$ across the eight cells: $0.089$ (FGLN, both
initialisations), $0.109$--$0.113$ (deep linear), $0.224$--$0.251$ (ReLU), $0.375$--$0.438$
(CReLU).

\paragraph{Conclusions.}
Only $9$--$44\%$ of what the loss asks of the operator is achievable by any weight update. The
ordering is explained by the rank of the step map. Measured at $L=6$, width $8$, the nullity of
$M$ is $19\%$ for CReLU, $53\%$ for ReLU, $83\%$ for deep linear and $93\%$ for FGLN: where the
gates are input-independent every layer moves the \emph{same} shared operator, so
$\mathrm{rank}(M) \le n_L n_0$ however many parameters the model has, and the fixed gates of FGLN
reduce it further. Reachability is therefore highest for the most expressive gating, which is the
opposite of the naive expectation that nonlinearity should obstruct the operator step.

\paragraph{A solver failure, and its resolution.}
Initially three of eight cells returned $\alpha > 1$ --- impossible --- with $\mathrm{ne}$ of
$28$, $54$ and $552$ instead of ${\sim}10^{-4}$. The cause was the failure mode described in
Section~\ref{sec:s2}: the undamped iterate drifts into $\ker M$, and it does so most where the
null space is largest, which is precisely the deep linear and FGLN cells. Tikhonov damping, for
which $\ker$ is trivial, together with float64 diagnostics, removed it: all $44$ probes now satisfy
$\alpha \le 1$ with residuals ${\sim}10^{-10}$. We record this because it is the reverse of the
natural suspicion --- the solve is least reliable on the \emph{simplest} models, and the
architectures the rest of Part II rests on are the best conditioned of the four.

\subsubsection{E4 --- What each optimiser does with the reachable part}

\paragraph{Motivation.}
This is the measurement the framework exists to make, and the one that endpoint studies cannot
make. Comparing optimisers by where they end confounds the quality of a step with the fact that
the arms have drifted to different points. Scoring every candidate from the \emph{same} weights
removes that confound entirely.

\paragraph{Setup.}
As in E3. A single optimiser (Adam, $\mathrm{lr} = 3\cdot10^{-3}$) drives the trajectory; at each
probe we stop and, from those weights and that batch, form the step each candidate would take.
Candidates: gradient descent; Adam; and the reference truncated at $k \in \{1,5,20,50\}$. Adam is
stateful, so its moments are carried along the driving trajectory at every step and ``Adam'' means
what Adam would do having followed this path --- which is exactly what it does, since Adam drives.
The probe at step $0$ is excluded because Adam's step is then identically zero and its cosine
undefined.

\paragraph{Measurement.}
$\cos_{\mathrm{op}}$ and $\cos_{\mathrm{w}}$ of Definition~\ref{def:ref} and the two norm ratios,
for every candidate, at every probe. Both cosines are scale-invariant and therefore carry no
learning rate. The table below reports $\cos_{\mathrm{op}}$; the weight-space comparison is
reported in E7, where the arms run freely and the two can be read against each other.

\paragraph{Status.} \textsc{done}, one seed, $8/8$ cells valid.

\paragraph{Results.}

\begin{center}\small
\begin{tabular}{llcccccc}
\toprule
arch & init & $\alpha$ & gd & Adam & $k{=}5$ & $k{=}20$ & $k{=}50$ \\
\midrule
CReLU  & looks-linear & 0.438 & 0.358 & $-0.070$ & 0.555 & 0.745 & 0.871 \\
CReLU  & Xavier       & 0.375 & 0.284 & $-0.087$ & 0.508 & 0.752 & 0.876 \\
FGLN   & orthogonal   & 0.089 & 0.444 & $\phantom{-}0.010$ & 0.808 & 0.927 & 0.938 \\
FGLN   & Xavier       & 0.089 & 0.412 & $\phantom{-}0.007$ & 0.769 & 0.875 & 0.896 \\
linear & orthogonal   & 0.109 & 0.583 & $-0.032$ & 0.896 & 0.994 & 0.999 \\
linear & Xavier       & 0.113 & 0.515 & $-0.025$ & 0.815 & 0.943 & 0.964 \\
ReLU   & orthogonal   & 0.251 & 0.410 & $-0.239$ & 0.582 & 0.815 & 0.913 \\
ReLU   & Xavier       & 0.224 & 0.323 & $-0.127$ & 0.530 & 0.799 & 0.903 \\
\midrule
\multicolumn{2}{l}{mean} & & 0.416 & $-0.070$ & 0.683 & 0.856 & 0.920 \\
\bottomrule
\end{tabular}
\end{center}

\paragraph{Conclusions.}
Three statements hold across every architecture and initialisation, and none of them involves a
learning rate.

\begin{enumerate}[leftmargin=*,itemsep=2pt,topsep=3pt]
  \item \textbf{Gradient descent is partially aligned}: it realises about $0.42$ of the reachable
        ideal direction on average, positive in $8/8$ cells and never close to $1$. It is moving
        the operator the way the loss asked, but inefficiently.
  \item \textbf{Adam is not aligned at all}: its mean $\cos_{\mathrm{op}}$ is negative in $6/8$
        cells and approximately zero in the other two. On average its step does not move the
        operator toward what the loss asked of it.
  \item \textbf{The truncation parameter recovers the ideal}: $0.68 \to 0.86 \to 0.92$ for
        $k = 5, 20, 50$, reaching $0.999$ for the deep linear orthogonal cell, where the reference
        step should be essentially exact and is.
\end{enumerate}

The consequence is the organising observation of this paper. \textbf{Adam trains these models at
least as well as gradient descent while doing something that is, at the operator level, close to
orthogonal to what the loss requests.} Two mechanisms must therefore be distinguished: motion of
the operator, and motion of everything else the parameterisation contains --- in a gated network,
the gate configuration. E5--E7 are designed to separate them.

\paragraph{Limitations.} One seed, one driving optimiser, one task, one depth. Whether the
conclusion depends on which optimiser drives the trajectory is not yet tested; see E7.

\subsubsection{E5 --- A reference configuration in which the task is actually learned}

\paragraph{Motivation.}
An endpoint comparison between two optimisers is meaningful only if both have learned the task. If
neither has, the comparison is between two failures and no conclusion about implicit bias survives
it. This experiment establishes the operating point, and exists because an earlier version of E6
did not have one.

\paragraph{Setup.}
MNIST-1D, $4000$ training, $1000$ validation, $1000$ test points, cross-entropy, batch $100$,
$12000$ steps. Bias-free networks, as \eqref{eq:fx} requires. Grid: architecture
$\in \{$CReLU, ReLU$\}$, depth $\in \{4, 8\}$, width $\in \{32, 64, 128\}$, Adam learning rate
$\in \{3\cdot10^{-3}, 10^{-2}, 3\cdot10^{-2}\}$; $36$ cells. Selection is on validation accuracy
and the test accuracy is reported \emph{at the selected step}, not as a maximum over the
trajectory.

\paragraph{Status.} \textsc{done} for the Adam arm; \textsc{todo} for the operator arm.

\paragraph{Results.}
Width is the binding constraint. At width $32$ the best test accuracy is $0.594$; at width $128$ it
is $0.695$. The best cell by validation is CReLU, $L=4$, width $128$, $\mathrm{lr} = 10^{-2}$, at
test $0.688$; the best cell that \emph{also} fits the training set is CReLU, $L=8$, width $128$,
$\mathrm{lr} = 3\cdot10^{-3}$, at test $0.653$ with training loss $0.055$. For ReLU the
corresponding cell is $L=8$, width $128$, $\mathrm{lr} = 10^{-2}$, at test $0.660$ with training
loss $0.05$. For reference, logistic regression on MNIST-1D reaches ${\sim}0.32$ and a
one-hidden-layer MLP ${\sim}0.68$.

\paragraph{Conclusions.}
The bias-free constraint costs little once the network is wide enough: at width $128$ these models
reach the accuracy expected of an MLP on this task. The configuration $L = 8$, width $128$ is
adopted as the reference operating point for E6, being the deepest setting in which both fitting
and generalisation are achieved.

This experiment also invalidates an earlier endpoint comparison conducted at width $32$, where the
best accuracy reached was $0.66$ and every configuration that drove the training loss to $10^{-4}$
reached only $0.44$--$0.57$. Worse, the single configuration that generalised acceptably never
reached the matched training loss and was therefore excluded from the comparison by construction.
Any rank gap measured there describes two over-fitted runs.

\subsubsection{E6 --- Endpoints at a configuration that learns}

\paragraph{Motivation.}
Given that Adam's step is unaligned at the operator level (E4) yet trains well (E5), the question
is whether the per-step difference accumulates into a different solution, and in what respect.

\paragraph{Setup.}
MNIST-1D at the reference configuration of E5: $L = 8$, width $128$, batch $100$, CReLU and ReLU,
$3$ seeds. Adam over a learning-rate grid spanning at least two orders of magnitude, against the
reference step at $k = 50$ and $\eta \in \{0.1, 0.3, 0.5\}$. Arms are compared \emph{at matched
training loss} rather than at a matched step count, which would measure optimisation speed instead.

\paragraph{Measurements.}
Test accuracy and test loss at matched training loss; and the effective rank of the operator, taken
as the participation ratio of its singular values averaged over inputs,
\begin{equation}
  \mathrm{PR}(x) \;=\; \frac{\big(\sum_i \sigma_i(P(x))\big)^2}{\sum_i \sigma_i(P(x))^2},
  \qquad
  \mathrm{PR} \;=\; \mathbb{E}_x\big[\mathrm{PR}(x)\big],
  \label{eq:pr}
\end{equation}
excluding inputs with $P(x) = 0$ and reporting their fraction separately, since such inputs carry
no spectrum and averaging their zero in drags the mean below the value $\mathrm{PR}$ takes on a
rank-one operator.

\paragraph{Status.} \textsc{todo} at the reference configuration. A previous run at width $32$ is
superseded for the reason given in E5.

\paragraph{Results.} \emph{[To be inserted.]}
At width $32$ the operator arm reached an effective rank higher by $+1.8$ to $+2.3$ in all six
cells, with Adam's rank flat across a five-point learning-rate grid, so the gap was not a tuning
artefact. Whether it survives at a configuration that learns the task is open.

\paragraph{Protocol.}
Two failure modes have already occurred in this project and are guarded against explicitly. A
best-over-grid statistic must not be read before every cell of the grid has completed; and the grid
must \emph{bracket} each method's optimum, since a method whose best value sits at the edge of its
grid has not been tuned. We report the full grid rather than the selected point, and seed-level
values rather than pooled means.

\subsubsection{E7 --- What Adam is doing instead: gates, dynamics and overfitting}

\paragraph{Motivation.}
E4 establishes that Adam's step is close to orthogonal to the operator request, and E5 that it
nonetheless learns the task. Something other than the operator is therefore being optimised. In a
gated network the natural candidate is the gate configuration: the operator $P(x)$ is what the
network computes \emph{given} a gate pattern, while the pattern itself is the combinatorial part of
the model and is invisible to $\Delta_\star$, which asks only that $P$ move.

\paragraph{Setup.}
The reference configuration of E5, with gradient descent, Adam and the reference step each running
freely, $3$ seeds, both architectures.

\paragraph{Measurements.} At each probe:
\begin{enumerate}[leftmargin=*,itemsep=2pt,topsep=3pt]
  \item \textbf{Gate occupancy.} Density $\rho = \mathbb{E}_{x,\ell,i}\big[\mathbf{1}[z_{\ell,i}(x) > 0]\big]$,
        and the dead-unit fraction, the proportion of hidden units inactive for every input.
  \item \textbf{Pattern diversity.} The mean normalised Hamming distance between the gate patterns
        of two independently drawn inputs, which measures how much of the network's behaviour is
        input-dependent at all.
  \item \textbf{Pattern churn.} The fraction of gates that flip between consecutive probes,
        separating motion of the gates from motion of the weights within a fixed pattern.
  \item \textbf{Operator rank.} \eqref{eq:pr}, as in E6.
  \item \textbf{Dynamics.} Training and test loss and accuracy against step, giving convergence
        speed and the point at which test loss begins to rise.
\end{enumerate}

\paragraph{Status.} \textsc{todo}.

\paragraph{Hypothesis under test.}
If Adam's advantage is gate reorganisation rather than operator motion, then relative to gradient
descent it should show higher pattern churn and greater pattern diversity while scoring lower on
$\cos_{\mathrm{op}}$; and the reference step, which optimises the operator directly, should show
the opposite. A null result --- indistinguishable gate statistics --- would mean the alignment gap
of E4 does not translate into any observable difference in what is learned, which would bound the
significance of E4 and must be reported as such.

% =========================================================================
\section{Conclusion}
\label{sec:conclusion}

Implicit bias is usually studied through the solutions an optimiser reaches. We have argued that it
can be studied through the steps it takes, provided one is willing to compute, at each point, the
weight update that would have realised what the loss asked of the operator. That reference step is
available in closed form for deep linear chains, where it is exact even at parameter values that
stop every derivative-based method, and as a least-squares solve for gated networks, where the
request is no longer fully realisable and the shortfall is itself a measurable property of the
architecture.

The resulting decomposition separates two things that endpoint studies necessarily conflate: what
the architecture permits, and what the optimiser does within what is permitted. The first is
optimiser-independent; the second is scale-invariant, and therefore free of the tuning objections
that comparisons between optimisers usually attract.

\emph{[Summary of findings to be inserted once the measurements of
Sections~\ref{sec:linear-exp} and \ref{sec:nonlinear-exp} are complete.]}

\paragraph{Limitations.}
The reference step is defined by a first-order model of how weights move the operator, and is
therefore meaningful only for steps small enough to stay in that regime; we report the step sizes
at which this holds. The least-squares solve becomes ill-conditioned with depth, and the
weight-space diagnostic $\cos_{\mathrm{w}}$ in particular is reliable only where the solve attains
the minimum-norm solution, which we verify per measurement rather than assume. Finally,
Remark~\ref{rem:fidelity} bears repeating: the diagnostics measure fidelity to the operator
gradient, which is not the same as optimisation quality, and nothing here should be read as a claim
that the reference step is a better training algorithm.

\newpage
\subsection*{AI use statement}

(This section is \textbf{required} and does not count toward the page limit.)

In this work, we used generative AI tools for [tasks with required disclosure].
We have not used generative AI tools for [other tasks with required disclosure],
and [the rest of the required disclosure tasks] are not applicable to this work.
Additionally, we used generative AI tools for [tasks with recommended
disclosure]. We have reviewed all AI-assisted work. [Elaborate. For example, “we
checked LLM-generated research ideas for potential plagiarism through a manual
literature survey”, “LLM-generated code was verified and tested for correctness
by 2 authors”, etc.]. We take responsibility for the final content of this work,
including text, claims or artifacts produced with the aid of generative AI.

See the ICLR 2027 AI Policy for Authors for more details. This statement should
not be more than 1 page.

\subsection*{Ethics statement}

\subsection*{Reproducibility statement}

\subsubsection*{Author Contributions}

\subsubsection*{Acknowledgments}

\bibliography{references}
\bibliographystyle{iclr2027_conference}

\appendix
\section{Appendix}

\section{Proofs}
\label{app:proofs}

\emph{[Proof of Proposition~\ref{prop:minimal}: Bellman recursion over depth with value function
$V_a(Y) = a\|Y\|_{S_{2/a}}^{2/a}$, and the arithmetic mean--geometric mean inequality for the
per-mode split. Proof of Proposition~\ref{prop:degen}: homogeneity and null-consistency of the
recursive construction. Proof of Proposition~\ref{prop:conserve}: the KKT conditions of
\eqref{eq:ref} give the shared multiplier \eqref{eq:kkt}, from which
$\Delta W_{\ell+1}^\top W_{\ell+1}$ and $\Delta W_\ell W_\ell^\top$ are transposes.]}

\section{Numerical behaviour of the least-squares solve}
\label{app:numerics}

\emph{[Divergence of LSQR past convergence on these step maps; failure of the two natural stopping
tests; the normal-equation stopping rule and the attained residuals by depth and architecture.]}

\section{Alternating least squares is not a minimal realisation}
\label{app:als}

\emph{[Measured ratio between the alternating-least-squares realisation and the minimal one, and
its consequences.]}

\end{document}
